\def\velkostpisma{10}
\def\velkostriadku{12.5}
\input makra.tex % nacitanie Ivanom pripravenych nastaveni a prikazov
\hyphenation{star-šov-stvo} % rozdelenie slov na konci riadku, treba tu uviest slova, ktore sam nepozna

\spravodaj{4}{2026}


\clanok{{\it ČO NÁS ČAKÁ V~APRÍLI}}


\clanok{Verš na mesiac}
V~apríli sa budeme učiť nový verš, ktorý dostala pre rok~2026 naša mládež:

„Cti Hospodina darmi svojho imania a prvotinami z~každého svojho výnosu. Tak sa ti naplnia stodoly hojnosťou a muštom pretečú tvoje kade.“ (Pr~3,9-10)


\clanok{Sedem slov}
Srdečne vás pozývame na Veľkopiatočné pašie do Nového evanjelického kostola na Legionárskej ulici v~Bratislave, 3.~4.~2026 o~18.00~hod., kde zaznie
nádherné oratórium francúzskeho skladateľa Théodore Dubois „Sedem slov Krista na kríži“.
Toto výnimočné dielo bude v~Bratislave uvedené po prvýkrát v~originálnej verzii v~podaní sólistov SND, speváckeho zboru Bratislava, spevokolu Meiningen (Nemecko) a za sprievodu symfonického orchestra Hilaris Chamber Orchestra pod vedením dirigenta Sebastiana Fuhrmanna (Nemecko).
Príďte a nechajte sa vniesť do hlbokej atmosféry Veľkého piatku a prežiť silný hudobný a duchovný zážitok, ktorý nám pripomína obeť nášho Spasiteľa Ježiša Krista a vieru vo večný život.

Tešíme sa na vašu účasť. Vstupné je dobrovoľné, príspevok na podporu realizácie projektu je vítaný.


\clanok{Spoločný obed v~SED}
V~apríli máme naplánovaný druhý spoločný obed. Stretneme sa v~nedeľu 5.~4. o~12.00~hod. v~priestoroch SED na Partizánskej. Prihlasovať sa nemusíte. Prineste si svoj obed a strávte tento čas pri spoločnom stole a rozhovoroch.


\clanok{Rozhovory na Palisádach}
V~sobotu dňa 11.~4.~2026 o~9.30~hod. na Palisádach prebehne ďalší z~„Rozhovorov na Palisádach“. Rečníkom bude Dávid M.~Chuchút s~témou „Liberalizmus, individualizmu, sloboda a kresťanstvo“.
Pozvánka s~viac informáciami je aj vo vašich e-mailových schránkach. Ak vás táto téma zaujíma, ste srdečne vítaní. Pozvite aj svojich priateľov.


\clanok{Jarná víkendovka}
Po ročnej pauze opäť plánujeme zborovú víkendovku v~penzióne Berea, termín je 17.~--~19.~4.~2026. Chceme spolu stráviť hodnotný neformálny čas a prehĺbiť vzájomné spoločenstvo. Témou je „SPOLU -- PRACOVAŤ -- ŽIŤ -- UCTIEVAŤ“, prevedú nás ňou bratia Janko Szőllős, Filip Barkóczi a Dávid M.~Chuchút. Podrobnejší program zverejníme čoskoro. Prihlasovací formulár sme rozosielali e-mailom. Veríme, že opäť spolu strávime hodnotný čas v~duchovnom spoločenstve ako aj vonku v~jarnej prírode.


\clanok{Sestry}
Sestry sa v~apríli nestretnú. Najbližšie stretnutie je v~máji a to v~stredu 20.~5.~2026.
\vfill\break


\clanok{{\it OHLIADNUTIE SA ZA MARCOM}}


\clanok{Správy zo staršovstva}
Drahí bratia a sestry,

staršovstvo zboru BJB Palisády sa v~marci stretlo dva razy a to 10.~3. a 24.~3.~2026. Na prvom stretnutí sa diskutovalo o~službe pre bezdomovcov so sestrou S.~Volentičovou a s~bratom T.~Valchářom o~spolupráci s~Connectom a správe Súľovskej od 1.~4.~2026. Zhodnotilo sa VZČZ a priebeh 1.~kola volieb, preberal sa denný tábor detí, podpora Dobrodiny, informovalo sa o~prípravách rekonštrukcie fasády a prebrali sa termíny DKDZ a KDZ. V~rámci rôzneho sa plánovala víkendovka v~Harmónii, výjazdová víkendovka staršovstva, spolupráca s~ukrajinským zborom v~Černihive a uvažovali o~spoločnom zhromaždení s~ukrajinským zborom.

Na druhom stretnutí staršovstva sme s~bratom Ľ.~Kešjarom diskutovali o~lepšej komunikácii medzi zložkami (technická podpora). Zhodnotili sme ZČZ, zvolenie kazateľa a samotný priebeh volieb. Ďalej sa schvaľovala nová nájomná zmluva k~Súľovskej~2. Plánoval sa program víkendovky. Boli sme informovaní o~začatí prípravných prác na rekonštrukcii fasády. Preberali sme kandidátku do volieb rady BJB. Brat kazateľ J.~Szőllős informoval o~komunikácii s~jednotlivcami s~pozastaveným členstvom. V~rôznom sa dohodol termín výjazdového staršovstva na Chvojnici 30.~--~31.~5.~2026. Ďalej spoločné bohoslužby s~ukrajinským zborom, nový organ, pôstne zamyslenia, rozhovory na Palisádach a cesta do Černihiva (detský tábor) 6.~--~11.~7.~2026.
Ďalšie stretnutie staršovstva sa uskutoční 7.~4.~2026.
\autor{Marcel Maďar}


\clanok{Voľby a čo ďalej?}
Bratia a sestry, ďakujeme nášmu Otcovi v~nebesiach za to, že nám v~Dávidovi M.~Chuchútovi daroval ďalšieho svojho služobníka. Hľadali sme Jeho vôľu pre náš zbor a aj keď sa to zdalo hektické a z~časového hľadiska veľmi natlačené do krátkeho obdobia, sme vďační za Jeho vedenie a veríme, že to, čo má On pre nás pripravené, je dobré a užitočné.
Môžu vzniknúť otázky, aké budú ďalšie kroky a čo nás kedy čaká?

Ďalšie kroky sú ovplyvnené viacerými faktormi. Jedným z~nich je aj voľba nových členov Rady BJB. Preto kroky, ktoré nás čakajú, budeme koordinovať s~novozvolenými členmi Rady.
Sú tu ale aj termíny dohodnuté v~minulosti, ktoré nemáme v~pláne meniť. Brat Dávid M.~Chuchút má naplánovanú prax v~zbore do konca júna~2026. Prax považujeme za dôležitú súčasť jeho prípravy na službu kazateľa v~zbore. To znamená, že Dávid M.~Chuchút svoju prax ukončí tak, ako to bolo dohodnuté pred jej začiatkom.

Od júla~2026 začne pre Dávida obdobie služby kazateľa nášho zboru a teda aj jeho volebné obdobie. Služba v~júli a auguste bude ovplyvnená zborovými tábormi a plánovanou dovolenkou koncom augusta.
V~septembri plánujeme Dávidovu inštaláciu za kazateľa zboru. Presný termín dnes nevieme určiť z~viacerých dôvodov. Ovplyvňuje ho voľba nových členov Rady BJB a dohoda termínu s~nimi. Ďalším faktorom je snaha o~to, aby na inštaláciu mohli prísť aj ľudia, ktorí sú Dávidovi blízki. Napriek týmto premenným máme snahu o~to, aby inštalácia mohla byť začiatkom septembra~2026.
Tak ako sme to robili v~minulosti, aj pri tejto inštalácii uvažujeme nielen so slávnostným zhromaždením, ale aj so spoločným obedom. Chceme stráviť spolu čas aj neformálne pri jedle, káve, koláčoch a hlavne pri rozhovoroch.

Prosíme vás, modlite sa aj za tento proces uvedenia Dávida M.~Chuchúta do služby kazateľa nášho zboru.
S~láskou a prianím Božieho vedenia a požehnania,

\autor{za staršovstvo zboru Peter Pribula}


\clanok{Pôstne zamyslenia}
Pôstne zamyslenia ma oslovili už samotným názvom, lebo pôst je pre mňa dlhodobo veľkou výzvou a zamyslenia dlhodobo milujem \smiley. Rada preto túto príležitosť netradičných bohoslužieb využívam a teší ma prítomnosť bratov a sestier, keď môžeme spolu v~tichu a prítmí počúvať zvestované Božie Slovo a premýšľať o~ňom ihneď a nie až doma, kde si na to už často nenájdeme čas.
\autor{Mirka Hovorková}


\clanok{„Rozhovory na Palisádach“ -- očami Slávky}
Asi by bolo lepšie, keby Izrael ostal v~mojej mysli hlavne ako spomienka na biblické postavy v~dlhých, pre nás netradičných odevoch s~pokrytou hlavou a príbehy na flanelografe v~besiedke.
Asi by bolo lepšie, keby sa Izrael spájal hlavne s~lahodným falafelom, kibucmi, ktoré rozkvitajú uprostred púšte, s~krajinou, ktorej dĺžka je 424~km a šírka na najužšom mieste iba 15~km.
Žijem však vo svete, ktorý hovorí o~Izraeli inak, a z~mnohých strán sa na nás chrlia úplne protichodné informácie.
Počúvame o~národe, ktorý sa bráni. O~národe, ktorý útočí. O~národe, ktorý je vyvolený. O~národe, ktorý je bohabojný. O~národe, ktorý je bezbožný.

Tak sme mnohí s~vdačnosťou počúvali slová Petra Šveca, ktorý sa nám snažil objasniť históriu, vznik štátu Izrael ako aj súčasnosť tejto geograficky drobnej a historicky bohatej krajiny. O~národe praotca Abraháma, kráľa Dávida, siroty Ester a mnohých iných.

Vďaka, Hospodine, za Tvoje Slovo, ktoré nám svieti na cestu, keď sa strácame v~lavíne mnohých cudzích slov...
\autor{Slávka Kráľová}
\vfill\break


\clanok{Zbierky v~marci}
\table{lr}{
Na misiu			& 413~€ \cr
Na investičný fond	& 375~€ \cr
}
\vskip1em

Pripomíname, že zbierky sa u~nás konajú pravidelne takto:
\vskip-1ex\begitems
* každú 2.~nedeľu v~mesiaci je zbierka venovaná misii a
* každú 4.~nedeľu je zbierka na investície.
\enditems

{\bf Pre dosporenie na rekonštrukciu} sme vypočítali príspevok na člena vo výške 150~€. Prispieť je možné formou zbierky počas 4.~nedieľ v~mesiaci alebo bankovým prevodom na zborový účet s~uvedením variabilného symbolu 777, prípadne aj osobne s.~Ľubke Kohútovej.

Bankové spojenie: SK36 0900 0000 0000 1147 1836, SWIFT: GIBASKBX

\n  3.	4.	Kirsten	PAYNTER;
\n  4.	4.	Viera	ŠKODÁK;
\n  6.	4.	Filip	BARKÓCZI;
\n  6.	4.	Jarmila	CIHOVÁ;
\n  6.	4.	Jana	ZAJACOVÁ;
\n 11.	4.	Daniel	MIKLETIČ;
\n 11.	4.	Erik	VAŠINA;
\n 16.	4.	Eva		MOLNÁROVÁ;
\n 19.	4.	Marta	PRIBULOVÁ;
\n 22.	4.	Alexander	Koloman	ERDÉLYI;
\n 28.	4.	Tanja	HAVROS;
\n 30.	4.	Ľuboš	DZURIAK;
\n 30.	4.	Jaroslav	VOLENTIČ;
\narodeniny


\program{
\p  1 ; st ;.;;.;;
\p  2 ; št ; 18.00 ; Pôstne zamyslenia 4. ;.;;
\p  3 ; pi ; 10.00 ; Pašie (D.~M.~Chuchút) ; 18.00 ; Sedem slov ;
\p  4 ; so ; 18.00 ; Mládež ;.;;
\p  5 ; ne ;  9.30 ; Bohoslužby (J.~Szőllős + VP) ; 12.00 ; Spoločný obed v~SED ;
\p  6 ; po ; 10.00 ; Výlet na Pajštún ;.;;
\p  7 ; ut ; 15.00 ; Biblická hodina pre seniorov (P.~Pivka) ;.;;
\p  8 ; st ;.;;.;;
\p  9 ; št ; 18.00 ; Biblická hodina (D.~M.~Chuchút) ;.;;
\p 10 ; pi ; 17.30 ; Dorast ;.;;
\p 11 ; so ; 18.00 ; Mládež ;.;;
\p 12 ; ne ;  9.30 ; Bohoslužby (R.~Nemec) ;.;;
\p 13 ; po ; 18.00 ; Modlitebná skupinka ;.;;
\p 14 ; ut ; 15.00 ; Biblická hodina pre seniorov (P.~Pivka) ;.;;
\p 15 ; st ;.;;.;;
\p 16 ; št ; 18.00 ; Biblická hodina (F.~Barkóczi) ;.;;
\p 17 ; pi ; 17.30 ; Dorast ;.;;
\p 18 ; so ; 18.00 ; Mládež ;.;;
\p 19 ; ne ;  9.30 ; Bohoslužby (D.~M.~Chuchút) ;.;;
\p 20 ; po ; 18.00 ; Modlitebná skupinka / skupinka „Rastieme v~Kristovi“ ;.;;
\p 21 ; ut ; 15.00 ; Biblická hodina pre seniorov (P.~Pivka) ;.;;
\p 22 ; st ;.;;.;;
\p 23 ; št ; 18.00 ; Biblická hodina (D.~M.~Chuchút) ;.;;
\p 24 ; pi ; 17.30 ; Dorast ;.;;
\p 25 ; so ; 18.00 ; Mládež ;.;;
\p 26 ; ne ;  9.30 ; Bohoslužby (F.~Barkóczi) ;.;;
\p 27 ; po ; 18.00 ; Modlitebná skupinka ;.;;
\p 28 ; ut ; 15.00 ; Biblická hodina pre seniorov (P.~Pivka) ;.;;
\p 29 ; st ;.;;.;;
\p 30 ; št ; 18.00 ; Biblická hodina (F.~Barkóczi) ;.;;
}

\tiraz
\bye
